% !TEX program = xelatex
\documentclass{misisarticle}
\usepackage{array}

\grnti{28.23.37, 73.47.12}
\asjc{1702}
\oecd{01.02.EP}

\title{Нейросетевая эвристика для ускорения поиска маршрута на дорожных сетях с зависящими от времени весами рёбер}

\misisauthor{Е.~В.~Мичков}{кафедра инженерной кибернетики НИТУ «МИСиС»}{Москва, Россия}{m2303197@edu.misis.ru}

\abstract{%
Рассматривается поиск самого быстрого маршрута на дорожной сети, где время проезда по дуге зависит от момента въезда на неё. Алгоритм TD-A* с нижней оценкой по весам свободного движения находит оптимальный маршрут, но в часы пик такая оценка сильно занижает оставшееся время и слабо сокращает область поиска. Предлагается эвристика, которая умножает нижнюю оценку на поправочный множитель, предсказанный многослойным перцептроном по координатам вершин и моменту времени. Параметр $\alpha \in [0, 1]$ задаёт долю поправки: при $\alpha = 0$ эвристика совпадает с допустимой нижней оценкой, при $\alpha = 1$ используется полное предсказание, которое может завышать оставшееся время и давать неоптимальный маршрут. Модель обучается на точных временах в пути, вычисленных алгоритмом TD-Dijkstra на графе Москвы и Московской области из OpenStreetMap с суточными профилями скорости. Метод сравнивается с TD-Dijkstra и TD-A* с ориентирами по числу раскрытых вершин, полному времени запроса с учётом расчёта нейросети и относительной ошибке времени в пути для разных рангов запроса и моментов отправления.
}

\keywords{%
поиск маршрута; дорожная сеть; зависящие от времени веса; алгоритм TD-A*; обучаемая эвристика; нейронная сеть
}

\begin{document}

\maketitle

\section{Введение}

Навигационный сервис строит маршрут с учётом момента выезда: одна и та же дуга дорожной сети в восемь утра и в полночь проезжается за разное время. Такую сеть описывают графом, в котором вес дуги задан функцией времени проезда от момента въезда на неё~\cite{kaufman1993}. Сведения о дорожной обстановке, из которых строят эти функции, дают в том числе бортовые камеры~\cite{pham2026} и лидары~\cite{vishnevsky2026}, а решения о движении на уровне отдельной полосы уже принимают нейросетевые агенты~\cite{kiselev2025}. На уровне сети остаётся задача выбора маршрута, и запрос к ней должен выполняться за миллисекунды.

Быстрее всего такие запросы обрабатывают методы с предобработкой графа: иерархии сжатия для зависящих от времени весов TCH~\cite{batz2013} и CATCHUp~\cite{strasser2021}. Цена скорости~--- жёсткость: иерархию нужно перестраивать или пересчитывать при изменении данных, а расширение постановки, например учёт запретов поворотов, требует изменения самой иерархии. Алгоритм A*~\cite{hart1968} принимает такие расширения почти без доработки~\cite{strasser2019}, но его скорость зависит от качества эвристики, то есть оценки оставшегося расстояния до цели.

Для зависящих от времени весов стандартная эвристика~--- нижняя оценка по минимальным весам дуг, которые соответствуют свободному движению~\cite{delling2007}. Она допустима, то есть не завышает оставшееся время, поэтому A* возвращает оптимальный маршрут. Однако в часы пик реальное время в пути заметно превышает время свободного движения, и оценка теряет точность именно тогда, когда запросов больше всего.

Обучаемая эвристика может учесть время суток и предсказать оставшееся время точнее. В работе Liashenko и Globa~\cite{liashenko2026} прямое предсказание времени в пути моделями LightGBM, GraphSAGE и GATv2 нарушало согласованность эвристики и увеличивало число раскрытых вершин по сравнению с классической эвристикой. Авторы получили выигрыш, когда сохранили допустимость: нейросеть предсказывает скорость, а эвристика вычисляется по геометрической формуле. Проверка проведена на графе из 2\,400 вершин на 10 парах вершин в три момента времени. Как ведёт себя недопустимая эвристика с контролируемой ошибкой на сети крупного города, эта работа не показывает.

Цель настоящей работы~--- оценить, насколько нейросетевая поправка к нижней оценке ускоряет TD-A* на дорожной сети Москвы и Московской области и какой ошибкой маршрута оплачивается это ускорение. Вклад работы:
\begin{itemize}
  \item эвристика в виде нижней оценки, умноженной на предсказанный поправочный множитель, с параметром $\alpha$, который непрерывно переводит её от допустимой к полностью предсказанной;
  \item схема получения обучающей выборки, в которой один запуск TD-Dijkstra даёт точные времена в пути до всех вершин графа;
  \item протокол сравнения, где время запроса включает расчёт нейросети, а ошибка маршрута измеряется отдельно для коротких и длинных запросов и для разных моментов отправления.
\end{itemize}

\section{Анализ работ}

Работы, близкие к теме, удобно сравнить по трём признакам: откуда берётся эвристика, сохраняет ли она оптимальность и как оценивалось ускорение. Сводка приведена в табл.~\ref{tab:related}.

\begin{table*}[htbp]
  \caption{Подходы к эвристике для поиска пути}
  \label{tab:related}
  \centering
  \begin{tabular}{>{\raggedright\arraybackslash}p{3.1cm}>{\raggedright\arraybackslash}p{4.4cm}>{\raggedright\arraybackslash}p{3.0cm}>{\raggedright\arraybackslash}p{5.6cm}}
    \toprule
    \textbf{Работа} & \textbf{Источник эвристики} & \textbf{Оптимальность} & \textbf{Оценка ускорения} \\
    \midrule
    TD-ALT~\cite{delling2007} & ориентиры по минимальным весам дуг & сохраняется & дорожные сети, зависящие от времени веса \\
    Двунаправленный TD-A*~\cite{nannicini2012} & ALT и обратный поиск по нижним оценкам & малая погрешность & дорожные сети, ускорение в несколько раз относительно Дейкстры \\
    CH-Potentials~\cite{strasser2019} & точные расстояния по нижним весам из CH & сохраняется & крупные дорожные сети, в том числе с прогнозами трафика \\
    PHIL~\cite{pandy2022} & графовая нейросеть, имитационное обучение & не гарантируется & разные графы, число раскрытых вершин \\
    CEA~\cite{futuhi2025} & нейросеть с функцией потерь против завышения & почти сохраняется & кубик Рубика \\
    TGNN~\cite{liashenko2026} & прогноз скорости и геометрическая формула & сохраняется & граф 2\,400 вершин, 10 пар вершин \\
    Настоящая работа & нижняя оценка и предсказанный множитель & регулируется $\alpha$ & граф Москвы и области, время запроса и ошибка маршрута \\
    \bottomrule
  \end{tabular}
\end{table*}

\subsection{Эвристики для зависящих от времени весов}

Поиск в сети с зависящими от времени весами сохраняет свойства алгоритма Дейкстры, если выполнено свойство FIFO: более поздний въезд на дугу не даёт более раннего прибытия~\cite{kaufman1993}. Delling и Wagner показали, что потенциалы ориентиров ALT остаются нижними оценками, если расстояния до ориентиров вычислены по минимальным весам дуг~\cite{delling2007}. Nannicini и соавторы использовали ту же идею для двунаправленного поиска: обратный поиск от цели идёт по нижним оценкам весов и ограничивает множество вершин для прямого поиска. Метод ускоряет запрос в несколько раз относительно алгоритма Дейкстры, если допустить небольшую погрешность маршрута~\cite{nannicini2012}. Strasser и Zeitz предложили CH-Potentials: иерархия сжатия, построенная по нижним весам, за короткое время даёт точное расстояние по этим весам и служит эвристикой для A*~\cite{strasser2019}. Во всех трёх подходах эвристика опирается на веса свободного движения и не зависит от момента отправления.

\subsection{Обучаемые эвристики}

Pándy и соавторы обучают графовую нейросеть PHIL на траекториях поиска и точных расстояниях. Её вычисление занимает постоянное время независимо от размера графа, а число раскрытых вершин на эталонных наборах сократилось в среднем на 58,5\,\% по сравнению с прежними методами~\cite{pandy2022}. Допустимость эвристики при этом не гарантируется. Futuhi и Sturtevant ставят обучение как задачу оптимизации с ограничением и вводят функцию потерь CEA, которая штрафует завышение оценки. На кубике Рубика она даёт почти допустимую эвристику, которая направляет поиск сильнее, чем сжатые базы шаблонов~\cite{futuhi2025}. Liashenko и Globa перенесли обучаемую эвристику на граф с зависящими от времени весами и пришли к выводу, что прямое предсказание времени в пути даёт несогласованную эвристику и лишние раскрытия вершин~\cite{liashenko2026}.

\subsection{Поиск с ограниченной неоптимальностью}

Pohl предложил взвешенный A*: эвристика умножается на коэффициент больше единицы, поиск раскрывает меньше вершин, а длина найденного пути превышает оптимальную не больше чем во столько же раз, если исходная эвристика допустима~\cite{pohl1970}. В настоящей работе роль коэффициента играет предсказанный множитель, который зависит от вершины и момента времени, поэтому такой гарантии нет и ошибку маршрута нужно измерять.

\subsection{Вывод из обзора}

Классические эвристики для зависящих от времени весов допустимы, но не учитывают момент отправления. Обучаемые эвристики учитывают данные, а в единственной найденной работе для зависящих от времени весов~\cite{liashenko2026} допустимость сохранена ценой формулы, которая ограничивает выигрыш, и проверка проведена на малом графе. Компромисс между ускорением и ошибкой маршрута для недопустимой обучаемой эвристики на сети крупного города остаётся неизмеренным. Этот пробел и закрывает настоящая работа.

\section{Постановка задачи}

Дорожная сеть задана ориентированным графом $G = (V, A)$. Каждой дуге $a \in A$ сопоставлена периодическая кусочно-линейная функция времени проезда $f_a(\tau)$ с периодом $\Pi = 24$~ч, где $\tau$~--- момент въезда на дугу. Для всех дуг выполнено свойство FIFO:
\begin{equation}
  \tau_1 \le \tau_2 \;\Rightarrow\; \tau_1 + f_a(\tau_1) \le \tau_2 + f_a(\tau_2).
  \label{eq:fifo}
\end{equation}

Запрос $q = (s, t, \tau_0)$ состоит из начальной вершины $s$, конечной вершины $t$ и момента отправления $\tau_0$. Для пути $P = (a_1, \ldots, a_k)$ моменты въезда определяются рекуррентно: $\tau_1 = \tau_0$, $\tau_{i+1} = \tau_i + f_{a_i}(\tau_i)$. Время в пути $T(P, \tau_0) = \tau_{k+1} - \tau_0$. Требуется найти путь из $s$ в $t$ с минимальным временем $T^*(s, t, \tau_0)$.

Алгоритм $\mathcal{A}$ возвращает путь $P_{\mathcal{A}}(q)$. Его качество на запросе описывают две величины: число раскрытых вершин $N_{\mathcal{A}}(q)$ вместе с временем выполнения запроса $t_{\mathcal{A}}(q)$, включая вычисление эвристики, и относительная ошибка
\begin{equation}
  \varepsilon_{\mathcal{A}}(q) = \frac{T(P_{\mathcal{A}}(q), \tau_0) - T^*(s, t, \tau_0)}{T^*(s, t, \tau_0)}.
  \label{eq:error}
\end{equation}

Эвристика $h(v, t, \tau)$ оценивает оставшееся время от вершины $v$ до $t$ при прибытии в $v$ в момент $\tau$. Она допустима, если $h(v, t, \tau) \le T^*(v, t, \tau)$ для всех $v$, $t$ и $\tau$. Для допустимой эвристики $\varepsilon = 0$.

Исследовательские вопросы:
\begin{enumerate}
  \item Во сколько раз нейросетевая эвристика сокращает $N$ и $t$ по сравнению с TD-Dijkstra и TD-A* с допустимой нижней оценкой?
  \item Как зависят ускорение и распределение $\varepsilon$ от параметра $\alpha$, ранга запроса и момента отправления?
  \item Окупается ли вычисление нейросети по полному времени запроса, а не только по числу раскрытых вершин?
\end{enumerate}

\section{Описание методов}

\subsection{Базовые алгоритмы}

TD-Dijkstra~\cite{kaufman1993} раскрывает вершины в порядке возрастания времени прибытия и при выполнении~\eqref{eq:fifo} находит оптимальный путь. Он служит эталоном для вычисления $T^*$ и нижней границей сравнения по скорости.

TD-A* раскрывает вершины в порядке возрастания ключа $\tau_v - \tau_0 + h(v, t, \tau_v)$, где $\tau_v$~--- текущее время прибытия в $v$. В качестве допустимой эвристики используется $h_{\mathrm{LB}}(v, t)$: потенциал ALT, вычисленный по минимальным весам дуг $\underline{w}_a = \min_\tau f_a(\tau)$ для набора из 16 ориентиров~\cite{delling2007}. Ориентиры выбираются заранее, расстояния от них и до них вычисляются один раз.

\subsection{Нейросетевая эвристика}

Предлагаемая эвристика не предсказывает время в пути целиком, а уточняет допустимую нижнюю оценку:
\begin{equation}
  h_\alpha(v, t, \tau) = h_{\mathrm{LB}}(v, t) \cdot \exp\bigl(\alpha \, g_\theta(x_{v,t,\tau})\bigr),
  \label{eq:heuristic}
\end{equation}
где $g_\theta$~--- многослойный перцептрон с параметрами $\theta$, а $\alpha \in [0, 1]$~--- доля поправки. При $\alpha = 0$ эвристика совпадает с $h_{\mathrm{LB}}$ и допустима. При $\alpha = 1$ она равна полному предсказанию и может завышать оставшееся время.

Такая форма выбрана по двум причинам. Во-первых, предсказание отношения к нижней оценке, а не самого времени, сужает диапазон целевой величины: перцептрону не нужно воспроизводить масштаб расстояний, который уже содержится в $h_{\mathrm{LB}}$. Во-вторых, параметр $\alpha$ позволяет получить всю кривую «ускорение~--- ошибка» одной обученной моделью, без переобучения.

Вектор признаков $x_{v,t,\tau}$ включает нормированные координаты вершин $v$ и $t$, расстояние по сфере между ними, значение $h_{\mathrm{LB}}(v, t)$ и момент $\tau$, закодированный парой $(\sin 2\pi\tau/\Pi, \cos 2\pi\tau/\Pi)$, чтобы полночь и 23:59 были близки. Перцептрон содержит два скрытых слоя по 64 нейрона с функцией активации ReLU. Модель такого размера не требует графического ускорителя, а время её вычисления входит в измеряемое время запроса.

\subsection{Обучающая выборка}

Целевая величина~--- $y = \ln\bigl(T^*(v, t, \tau) / h_{\mathrm{LB}}(v, t)\bigr)$, функция потерь~--- среднеквадратичная ошибка. Для разметки используется свойство TD-Dijkstra: запуск из вершины $v$ в момент $\tau$ даёт точные $T^*(v, u, \tau)$ сразу для всех вершин $u$. Поэтому выборка строится так:
\begin{enumerate}
  \item выбирается случайная вершина $v$ и случайный момент $\tau$ в пределах суток;
  \item TD-Dijkstra из $v$ в момент $\tau$ вычисляет времена до всех вершин;
  \item из полученных вершин случайно отбирается подмножество целей $t$ с равномерным распределением по рангу, чтобы короткие и длинные пары были представлены одинаково.
\end{enumerate}
Вершины для обучения и для проверки не пересекаются, чтобы оценивать обобщение на новые пары, а не запоминание.

\subsection{Поиск с недопустимой эвристикой}

При недопустимой эвристике вершина может быть раскрыта раньше, чем найден оптимальный путь до неё. Повторное раскрытие таких вершин восстанавливает часть точности, но в худшем случае резко увеличивает число раскрытий~--- именно этот эффект Liashenko и Globa наблюдали для прямого предсказания времени~\cite{liashenko2026}. В настоящей работе вершины повторно не раскрываются, поэтому число раскрытий не превышает $|V|$, а цена недопустимости проявляется только в ошибке~\eqref{eq:error}. Значение эвристики для вершины вычисляется один раз, при первом её достижении, и далее берётся из кэша запроса.

\subsection{Данные}

Граф строится из выгрузки OpenStreetMap по Москве и Московской области: оставляются дороги, доступные легковому автомобилю, и наибольшая сильно связная компонента. Суточные профили скорости формируются самостоятельно по классу дороги: скорость свободного движения снижается в утренний и вечерний пики, и функции времени проезда проверяются на свойство FIFO. Профили синтетические, поэтому выводы работы относятся к модели загрузки, а не к реальному трафику Москвы. Это ограничение указывается при интерпретации результатов.

\subsection{Протокол эксперимента}

Тестовые запросы выбираются по рангу Дейкстры: для случайной вершины $s$ и момента $\tau_0$ целью $t$ берётся вершина, раскрытая TD-Dijkstra под номером $2^r$, $r = 6, \ldots, \lfloor \log_2 |V| \rfloor$. Моменты отправления берутся из ночного интервала, утреннего пика и дневного интервала. Сравниваются TD-Dijkstra, TD-A* с $h_{\mathrm{LB}}$ и TD-A* с $h_\alpha$ для $\alpha \in \{0{,}25;\, 0{,}5;\, 0{,}75;\, 1\}$. Для каждого алгоритма регистрируются $N$, $t$ и $\varepsilon$; по $\varepsilon$ приводятся медиана, 95-й процентиль, максимум и доля запросов с $\varepsilon = 0$. Время запроса измеряется на одном компьютере при фиксированных условиях, все алгоритмы используют одно представление графа и одну очередь с приоритетами, чтобы разница во времени объяснялась эвристикой, а не реализацией.

\renewcommand{\refname}{ЛИТЕРАТУРА}
\begin{thebibliography}{99}

\bibitem{kaufman1993} Kaufman D.~E., Smith R.~L. Fastest paths in time-dependent networks for intelligent vehicle-highway systems application // IVHS Journal. 1993. Vol.~1, no.~1. P.~1--11. DOI: 10.1080/10248079308903779.

\bibitem{pham2026} Фам М.~Х. Определение дистанции до впереди идущих транспортных средств по видеопоследовательности для системы предупреждения столкновений с использованием YOLOv8 и Depth-Anything-V2 // Искусственный интеллект в промышленных, коммерческих, медицинских и финансовых приложениях: сборник статей научно-технического семинара. Вып.~6. М.: НИТУ «МИСИС», 2026. С.~117--124.

\bibitem{vishnevsky2026} Вишневский М.~Р., Фёдоров И.~Ю. Повышение пространственного разрешения механических LiDAR-систем с использованием виртуального зеркального канала и автоматической калибровки // Труды ИСП РАН. 2026. Т.~38, вып.~3, ч.~3. С.~101--114. DOI: 10.15514/ISPRAS-2026-38(3)-38.

\bibitem{kiselev2025} Киселев С.~Д., Алтунян А.~В. Локальные методы планирования траекторий на основе обучения с подкреплением // Искусственный интеллект в промышленных, коммерческих, медицинских и финансовых приложениях: сборник статей научно-технического семинара. Вып.~3. М.: НИТУ «МИСИС», 2025. С.~84--89.

\bibitem{batz2013} Batz G.~V., Geisberger R., Sanders P., Vetter C. Minimum time-dependent travel times with contraction hierarchies // ACM Journal of Experimental Algorithmics. 2013. Vol.~18. DOI: 10.1145/2444016.2444020.

\bibitem{strasser2021} Strasser B., Wagner D., Zeitz T. Space-efficient, fast and exact routing in time-dependent road networks // Algorithms. 2021. Vol.~14, no.~3. Art.~90. DOI: 10.3390/a14030090.

\bibitem{hart1968} Hart P.~E., Nilsson N.~J., Raphael B. A formal basis for the heuristic determination of minimum cost paths // IEEE Transactions on Systems Science and Cybernetics. 1968. Vol.~4, no.~2. P.~100--107. DOI: 10.1109/TSSC.1968.300136.

\bibitem{strasser2019} Strasser B., Zeitz T. A fast and tight heuristic for A* in road networks // arXiv preprint. 2019. arXiv:1910.12526.

\bibitem{delling2007} Delling D., Wagner D. Landmark-based routing in dynamic graphs // Experimental Algorithms (WEA 2007). Lecture Notes in Computer Science. Vol.~4525. Berlin: Springer, 2007. P.~52--65. DOI: 10.1007/978-3-540-72845-0\_5.

\bibitem{liashenko2026} Liashenko A., Globa L. Accelerating A* algorithm based search in time-dependent graphs with learned heuristics // Radioelectronic and Computer Systems. 2026. No.~1. P.~179--191. DOI: 10.32620/reks.2026.1.12.

\bibitem{nannicini2012} Nannicini G., Delling D., Schultes D., Liberti L. Bidirectional A* search on time-dependent road networks // Networks. 2012. Vol.~59, no.~2. P.~240--251. DOI: 10.1002/net.20438.

\bibitem{pandy2022} Pándy M., Qiu W., Corso G., Veličković P., Ying R., Leskovec J., Liò P. Learning graph search heuristics // Proceedings of the First Learning on Graphs Conference. PMLR. 2022. Vol.~198. P.~10:1--10:13.

\bibitem{futuhi2025} Futuhi E., Sturtevant N.~R. Learning admissible heuristics for A*: theory and practice // arXiv preprint. 2025. arXiv:2509.22626.

\bibitem{pohl1970} Pohl I. Heuristic search viewed as path finding in a graph // Artificial Intelligence. 1970. Vol.~1, no.~3--4. P.~193--204. DOI: 10.1016/0004-3702(70)90007-X.

\end{thebibliography}

\end{document}
